\section{Proposed UNLamb Benchmark}
\label{sec:benchmark}

To systematically investigate the influence of knowledge popularity on machine unlearning, we introduce \textbf{UNLamb}, a new benchmark constructed from the PopQA dataset \cite{mallen2023not}. We chose PopQA as our foundation due to its direct mapping of natural language questions to structured Wikidata facts, which crucially includes quantifiable popularity score for each entity.
To further motivate the distinction between popular and rare knowledge, we created the largest split—15\% rare and popular facts (which subsumes the 5\% and 10\% splits.) -- and evaluated the original LLaMA-3 models of various sizes on this split (See Table~\ref{tab:orig_forget15_by_size}); the results clearly reaffirm that the models retain popular knowledge substantially better than rare knowledge.


\begin{table}[h!t]
  \centering
  \small
  \setlength{\tabcolsep}{6pt}
  \begin{tabular}{|l|c|c|c|}
    \hline
     & \textbf{1B} & \textbf{3B} & \textbf{8B} \\ 
    \hline
    rare\_forget15 & 0.010 & 0.027 & 0.051 \\
    \hline
    popular\_forget15 & 0.099 & 0.381 & 0.536 \\
    \hline
  \end{tabular}
  % \caption{ROUGE-L Scores for LLaMA-3 Unlearning on Rare vs. Popular 15\% Datasets Across Model Sizes.}
  \caption{ROUGE-L scores evaluating original pretrained LLaMA-3 models of various sizes on 15\% rare and popular forget-set splits.}
  \label{tab:orig_forget15_by_size}
\end{table}






\subsection{Corpus Definition and Popularity Metric}

The UNLamb corpus $\mathcal{D}$ consists of $N=11,600$ factual triplets, where each fact $f_i$ is a tuple $(q_i, s_i, p_i, a_i)$:
\begin{itemize}
    \item $q_i \in \mathcal{Q}$ is a question in natural language (e.g., \enquote{What is the capital of France?}).
    \item $s_i \in \mathcal{E}$ is the subject entity from Wikidata (e.g., \enquote{France}).
    \item $p_i \in \mathcal{R}$ is the Wikidata property (e.g., \texttt{1082}, integer population).
    \item $a_i \in \mathcal{A}$ is the single, canonical answer (e.g., \enquote{Paris}).
\end{itemize}


For each fact $f_i$, its popularity score, $\operatorname{pop}(f_i)$, is defined as the average monthly Wikipedia page views for its subject entity $s_i$.
This metric serves as a strong proxy for a fact's prevalence, as Wikipedia is a core component of LLM pretraining corpora~\cite{gao2020pile, dodge2021documenting}. Our evaluation is therefore more realistic, targeting the removal of organically acquired knowledge rather than artifacts of fine-tuning.


By sorting the entire corpus $\mathcal{D}$ based on this score in ascending order, we obtain an ordered sequence of facts $f_{(1)}, f_{(2)}, \dots, f_{(N)}$, where $\operatorname{pop}(f_{(1)}) \le \operatorname{pop}(f_{(2)}) \le \dots \le \operatorname{pop}(f_{(N)})$.

\subsection{Popularity-Stratified Forget and Retain Sets}
\label{subsec:popular-strat}

The core of our benchmark design is to create unlearning tasks targeting the two extremes of the popularity spectrum. We select forgetting percentages \(p \in \{5\%,10\%,15\%\}\). This choice is inspired by benchmarks like TOFU \cite{maini2024tofu}, which uses splits of \(\{1\%,5\%,10\%\}\), but we shift the range upwards because a 1\% split provides too little data to produce stable, informative unlearning signals \cite{pal2025llm}. For a given percentage $p$, we define the number of samples to forget from each extreme as $k = \lfloor \frac{p}{100}\mathcal{D} \rfloor$. Given our full benchmark size \( \mathcal{D} = 11600 \), this corresponds to forget sets (\( F_p^{\text{rare}} \), \( F_p^{\text{pop}} \)) of 582 (\( 5\% \)), 1160 (\( 10\% \)), and 1750 (\( 15\% \)) samples, respectively.
This allows us to define two distinct \emph{forget sets}:
\begin{itemize}
    \item \textbf{Rare Forget Set ($F_p^{\text{rare}}$):} The $k$ least popular facts in the corpus.
    \begin{equation}
      F_p^{\text{rare}} = \{ f_{(1)}, \dots, f_{(k)} \}
    \end{equation}
    \item \textbf{Popular Forget Set ($F_p^{\text{pop}}$):} The $k$ most popular facts in the corpus.
    \begin{equation}
      F_p^{\text{pop}} = \{ f_{(D - k)}, \dots, f_{D} \}
    \end{equation}
\end{itemize}

The remaining facts form the \textbf{Retain Set} \(R_p\), which quantifies how much general knowledge remains preserved and untouched by the unlearning procedure.
\begin{equation}
  R_p = \mathcal{D} \setminus \bigl(F_p^{\text{rare}} \cup F_p^{\text{pop}}\bigr)
  \label{eq:retain-construction}
\end{equation}


\subsection{Probing Unlearning Locality: Validation Splits}
\label{subsec:duplicates_splits}

A successful unlearning algorithm should be precise, affecting only the targeted knowledge. To measure unintended side effects, or the \emph{locality} of forgetting, we construct validation splits from the retain set $R_p$. These splits contain facts that are related to the forget sets by either sharing the same subject or the same answer.

Let $S(F) = \{s_i \mid (q_i, s_i, p_i, a_i) \in F\}$ be the set of all subjects in a fact set $F$, and $A(F)$ be the set of all answers. We define the following validation sets:
\begin{itemize}
    \item \textbf{Duplicate-Subject Sets ($V_{\text{DS}}$):} These sets test if the model forgets other facts about a subject $s(f)$ when one fact about it is unlearned.
    \begin{align}
      V_{\text{DS}, p}^{\text{rare}} &= \{ f \in R_p \mid s(f) \in S(F_p^{\text{rare}}) \} \label{eq:vds-rare} \\
      V_{\text{DS}, p}^{\text{pop}}  &= \{ f \in R_p \mid s(f) \in S(F_p^{\text{pop}}) \} \label{eq:vds-pop}
    \end{align}
    
    \item \textbf{Duplicate-Answer Sets ($V_{\text{DA}}$):} These sets test if recalling an answer is impaired when one question leading to it is unlearned.
    \begin{align}
      V_{\text{DA}, p}^{\text{rare}} &= \{ f \in R_p \mid a(f) \in A(F_p^{\text{rare}}) \} \label{eq:vda-rare} \\
      V_{\text{DA}, p}^{\text{pop}}  &= \{ f \in R_p \mid a(f) \in A(F_p^{\text{pop}}) \} \label{eq:vda-pop}
    \end{align}
\end{itemize}

% Together, these splits enable a fine-grained evaluation of unlearning algorithms, measuring not only their efficacy on the target sets ($F_p^{\text{rare}}, F_p^{\text{pop}}$) but also their collateral impact on related knowledge ($V_{\text{DS}}, V_{\text{DA}}$) and general knowledge ($R_p$).

% The visualization of the splitting subsets for \(p=10\)\% is presented in Figure~\ref{fig:UNLamb}.

Together, these splits enable a fine-grained evaluation of unlearning algorithms, measuring not only their efficacy on the target sets ($F_p^{\text{rare}}, F_p^{\text{pop}}$) but also their collateral impact on related knowledge ($V_{\text{DS}}, V_{\text{DA}}$) and general knowledge ($R_p$). Figure~\ref{fig:UNLamb} summarizes the construction for \(p=10\%\): the top panel shows how the retain set \(R_p\) is derived from the full fact set \(D\) by excluding the forget splits \(F_p^{\text{rare}}\) and \(F_p^{\text{pop}}\) (Eq.~\eqref{eq:retain-construction}), the bottom left depicts the generation of duplicate-subject validation examples for rare and popular subsets (Eqs.~\eqref{eq:vds-rare}--\eqref{eq:vds-pop}), and the bottom right shows the creation of duplicate-answer examples (Eqs.~\eqref{eq:vda-rare}--\eqref{eq:vda-pop}).





% \textit{UNLamb is hosted on the Hugging Face Hub: \url{https://huggingface.co/datasets/SwetieePawsss/UNLamb}.}